\section{Conclusion}

This paper reframes LLM-as-a-Judge reliability as a question about the
evaluation dimensions considered before scoring. The hard-gold diagnostic
shows that a single-model evaluation-criteria policy leaves substantial
human-gold dimensions uncovered, motivating evaluation blind spots as a
first-class failure mode rather than a generic criteria-quality issue.

Blind-Spot Coverage makes this failure mode measurable: it rewards semantic
coverage of human evaluation dimensions while controlling redundancy and
invalid or hallucinated criteria. The resulting BSC reward gives a concrete
way to test open-ended criteria elicitation under verifiable-reward
optimization, despite lacking a single exact target string.

The empirical claims in this manuscript are intentionally evidence-gated. The
minimal motivation gate supports the hard-gold blind-spot diagnostic and its
robustness artifacts for Section 2 only; it is separate from the full real-run
C0--C14 readiness matrix. Statements about whether RLVR/GRPO warrants a
dimension-level coverage-change or recovery statement, whether a result is supported
by downstream utility, or whether semantic-space coverage evidence is
paper-facing are permitted only when the hard-gold, contamination, ablation,
downstream, dimension-transition, and visualization gates pass.
In particular, a BSC coverage change alone remains a metric result until redundancy,
hallucination, and held-out downstream utility evidence support the same
conclusion.
This discipline is central to the paper: the contribution is not a claim of a
better policy by construction, but an auditable framework for discovering,
measuring, and testing when dimension-level coverage-change or recovery wording
about evaluation blind spots is supported.
